\section{Problem Formulation}
\label{3}
\subsection{Problem Setting}
We introduce a novel problem setting called Continual Learning Unlearning (CLU) also known as Continual Adaptation or simply Adaptations, a framework for dynamically modifying model knowledge defined as learned class mappings encoded in parameters. This setting brings together two complementary goals: (i) selectively removing specific targeted knowledge (machine unlearning) and (ii) acquiring new targeted knowledge for existing pre-trained models (continual learning), all while maintaining performance on the remaining knowledge in a continual form. To develop a solution for this problem, we first present the most basic case, in which only one adaptation task is performed, and then expand this formulation to the continuous scenario, which includes a series of adaptation activities.

Let $\mathcal{M}$ be a model pre-trained on the dataset $D$. We regard $\mathcal{M}$ as a mapping function $f_\mathcal{M} : \mathcal{X}_D \rightarrow \mathcal{Y}_D$, where $\mathcal{X}_D$ and $\mathcal{Y}_D$ denote the input and output spaces associated with $D$, respectively. Our objective is to selectively discard certain knowledge while adding new knowledge and retaining the rest. We assume $D_f$ be the dataset to be forgotten, $D^\text{\textit{full}}_r$ the dataset to be retained, and $D_{\text{new}}$ the new dataset to be learned, satisfying $D = D^\text{\textit{full}}_r \cup D_f$, $D^\text{\textit{full}}_r \cap D_f = \emptyset$, and $D_{\text{new}} \cap D = \emptyset$.

In practice, storing the full retain set is prohibited under General Data Protection Regulation (GDPR) and California Consumer Privacy Act (CCPA) regulations for privacy-sensitive data, while for non-private data, retraining remains computationally expensive due to large-scale models and datasets. We therefore maintain a small replay buffer $D_r \subset D^\text{\textit{full}}_r$ such that $|D_r| \ll |D^\text{\textit{full}}_r|$ (at least 10\% of $|D^\text{\textit{full}}_r|$), here $D^\text{\textit{full}}_r$ represents full retaining data \footnote{From this point, $D_r$ denotes the experience replay buffer unless specified.}. The consequences of insufficient or absent buffer size are discussed in detail in Section~\ref{8.1}.


Before adaptation, $\mathcal{M}$ performs well on $D_r$ and $D_f$ but poorly on $D_{\text{new}}$, i.e.,
{\small
\begin{equation}
f_{\mathcal{M}} :
\mathcal{X}_{D_f} \rightarrow \mathcal{Y}_{D_f} \; ; \;
\mathcal{X}_{D_r} \rightarrow \mathcal{Y}_{D_r} \; ; \;
\mathcal{X}_{D_{\text{new}}} \cancel{\rightarrow} \mathcal{Y}_{D_{\text{new}}}.
\label{eq:1}
\end{equation}
}

Let the adaptation algorithm be $\mathcal{F}$ for more details please refer Section \ref{4} which modifies the model $\mathcal{M}$ to obtain $\mathcal{M'}$ by using $\mathcal{F}$ such that $\mathcal{M}' = \mathcal{F}(\mathcal{M}, D_f, D_r, D_{\text{new}})$ such that $\mathcal{M'}$ holds a new mapping relationship $f_{\mathcal{M}'}$ as 

{\small
\begin{equation}
f_{\mathcal{M'}} :
\mathcal{X}_{D_f} \cancel{\rightarrow} \mathcal{Y}_{D_f} \; ; \;
\mathcal{X}_{D_r} \rightarrow \mathcal{Y}_{D_r} \; ; \;
\mathcal{X}_{D_{\text{new}}} \rightarrow \mathcal{Y}_{D_{\text{new}}}.
\label{eq:2}
\end{equation}
}
Here, \( \cancel{\rightarrow} \) denotes that the mapping no longer holds (i.e., the model has forgotten the corresponding mapping), while \( \rightarrow \) indicates that the mapping still holds. 

We now extend this problem to the continual setting, where the model is required to sequentially adopt new knowledge, involving both unlearning and learning across tasks $T$. These tasks may be triggered by requests from the user, the owner, or both. Here, $t$ is an index representing the task number, ranging from $0, 1, 2, \ldots, t, \ldots, T$. let $D_r^\text{(t)}$ be the data to be retained, $D_{f}^\text{(t)}$ be the data to be forgotten and $D_\text{new}^\text{(t)}$ be the data to be learned at $t_\text{th}$ step respectively.

The algorithm $\mathcal{F}$ is applied as:
\[
    \mathcal{M}^\text{(t)} = \mathcal{F}(\mathcal{M^\text{{(t - 1)}}}, D_{f}^\text{(t)}, D_{r}^\text{(t)}, D_{\text{new}}^\text{(t)})
\]

The adaptation algorithm $\mathcal{F}$ takes the previous-step model $\mathcal{M^\text{{(t - 1})}}$ along with the datasets 
$D_{r}^\text{(t)}$, $D_{f}^\text{(t)}$, and $D_{\text{new}}^\text{(t)}$ to produce the updated model $\mathcal{M^\text{(t)}}$. Thus, $\mathcal{F}$ handles adaptation requests sequentially, starting from $\mathcal{M}$ and generating  a sequence of models $\mathcal{M^\text{(1)}}, \; \mathcal{M^\text{(2)}}, \; \mathcal{M^\text{(3)}}, \; \ldots, \; \mathcal{M^\text{(t)}}, \; \ldots, \; \mathcal{M^\text{(T)}}$,  where $\mathcal{M^\text{(t)}}$ represents the modified model after the $t$-th adaptation task.

After step $t$, the adapted model $\mathcal{M}^{(t)}$ must satisfy:

\begin{subequations}
\label{Eq:3}
\begin{align}
    f_{\mathcal{M}^{(t)}} &: \mathcal{X}_{D_{f}^{(i)}} \cancel{\rightarrow} \mathcal{Y}_{D_{f}^{(i)}}, \quad \forall \, i \leq t \label{Eq:3a} \\
    f_{\mathcal{M}^{(t)}} &: \mathcal{X}_{D_{r}^{(t)}} \rightarrow \mathcal{Y}_{D_{r}^{(t)}} \label{Eq:3b} \\
    f_{\mathcal{M}^{(t)}} &: \mathcal{X}_{D_{\text{new}}^{(i)}} \rightarrow \mathcal{Y}_{D_{\text{new}}^{(i)}} \label{Eq:3c}
\end{align}
\end{subequations}


We define successful CLU as 
\begin{itemize}
    \item Forgetting: $\text{Acc}(\mathcal{M}^{(t)}, D_f^{(i)}) \leq \frac{1}{C}, \quad \forall \, i \leq t$
    \item Retention: $\text{Acc}(\mathcal{M}^{(t)}, D_r^{(t)}) \approx \text{Acc}(\text{Oracle}, D_r^{(t)})$
    \item Learning: $\text{Acc}(\mathcal{M}^{(t)}, D_{\text{new}}^{(t)}) \approx \text{Acc}(\text{Oracle}, D_{\text{new}}^{(t)})$
\end{itemize}
where $C$ is the number of classes and Oracle denotes a model trained directly on target classes only.

\subsection{Assumptions}
\label{3.2}
BID-LoRA assumes the following conditions hold:
\begin{enumerate}
    \item A replay buffer $D_r$ with $|D_r| \geq 0.1|D^\text{\textit{full}}_r|$ is available.
    \item Learning and unlearning occur simultaneously; if only one is needed, the corresponding loss is masked (see Section~\ref{4.3}).
    \item Data partitions satisfy $D = D^\text{\textit{full}}_r \cup D_f$, $D^\text{\textit{full}}_r \cap D_f = \emptyset$, $D_{\text{new}} \cap D = \emptyset$.
    \item \textcolor{red}{All adapter matrices permanently deleted post-merge.}
\end{enumerate}


